This paper has introduced a graph-guided prompting workflow intended to improve the reliability of long-form scientific writing produced with large language models. The central components of the workflow are (1) a document graph that encodes section and subsection dependencies and that drives a leaf-first scheduling policy for generation, (2) an evidence retrieval and ranking protocol that selects candidate snippets for injection into section prompts, (3) a strict JSON schema and evidence-aware prompt templates that constrain section outputs and require explicit source citations for factual statements, and (4) an assembly procedure that converts validated JSON section artifacts into a LaTeX manuscript with an accompanying refs.bib file. Together, these artifacts are presented as a cohesive pipeline designed to make section generation, post-hoc audits, and iterative edits machine readable and reproducible within standard LaTeX/bib workflows.

A methodological emphasis of this work is conservative claim-making and mandatory citation of sources for factual assertions: the generation templates and the JSON schema explicitly require citation entries for non-opinion statements, and the validation steps prioritize detection of unsupported assertions before assembly. The evaluation protocol described in this manuscript focuses on citation fidelity and factual consistency as primary metrics; further methodological details and quantitative results are reported in Section 5. We intentionally avoid broad generalization beyond the scope of the datasets and retrieval configurations evaluated in order to prioritize transparency about the workflow's capabilities and limits.

The paper also documents practical limitations encountered in the current study, including sensitivity to retrieval quality, constraints imposed by small or curated dataset sizes, and challenges associated with transferring the workflow across domains. (General background knowledge.) These limitations motivate conservative interpretation of empirical findings and motivate the design choices emphasizing auditability and schema-enforced outputs (see Section 6 for a fuller discussion).

Future work will pursue larger-scale, multi-domain evaluations and the development of community-shared datasets and toolkits to facilitate reproducible benchmarking of citation-aware generation systems. (General background knowledge.) Additional extensions include automating citation–assertion consistency checks, strengthening provenance tracking for injected evidence, and integrating the pipeline with reproducible computational environments to streamline end-to-end audits and collaborative editing.

In closing, this work contributes a practical, audit-oriented approach to structuring long-form scientific generation around an explicit document graph, evidence-aware prompt engineering, and machine-readable section artifacts. By foregrounding citation discipline and validation at the section level, we aim to reduce unsupported assertions while preserving an editable, LaTeX-compatible output format. We encourage further evaluation and community collaboration to refine retrieval, evaluation, and scaling strategies that advance reliable, evidence-grounded long-form generation.